\section{Размерности $9$--$12$: объёмный экран и ламинирование}\label{sec:dim912}

Все конструкции предыдущих разделов ищут ядро внутри \emph{фиксированной}
родительской решётки, а затем деформируют её метрику. В размерности~$9$ этот
путь упирается в вычислительную стену: запрещённое множество решётки $A_9^*$
содержит $389\,387$ векторов, и локальный поиск в нём безрезультатен. Оказалось,
что дело не в оптимизаторе.

\subsection{Объёмный экран Минковского}\label{ssec:mink}

Условие правильности раскраски $D(v)\ge\diam V_0$ для всех $v\in\Gamma
\setminus\{0\}$ равносильно геометрическому утверждению
\[
\Gamma\cap\operatorname{int}K=\{0\},\qquad K=2\bigl(V_0\oplus R\,B\bigr),
\]
где $R$ --- радиус покрытия, $B$ --- единичный шар. Тело $K$ центрально
симметрично и выпукло как сумма Минковского симметричного многогранника и шара,
поэтому применима теорема Минковского о выпуклом теле.

\begin{proposition}[объёмная нижняя граница индекса]\label{prop:mink}
\textbf{[Т]} Для любой подрешётки $\Gamma\subset\Lambda$, задающей правильную
раскраску,
\begin{equation}\label{eq:mink}
k=[\Lambda:\Gamma]\ \ge\ \frac{\vol\bigl(V_0\oplus R\,B\bigr)}{\det\Lambda}\,,
\end{equation}
и, по неравенству Брунна--Минковского, в замкнутой форме
$k\ge\bigl(1+R\,\omega_n^{1/n}/(\det\Lambda)^{1/n}\bigr)^n$, где $\omega_n$ ---
объём единичного шара.
\end{proposition}

\begin{proof}
Из $\Gamma\cap\operatorname{int}K=\{0\}$ по теореме Минковского
$\vol K\le2^n\det\Gamma$; при этом $\vol K=2^n\vol(V_0\oplus RB)$, откуда
$\det\Gamma\ge\vol(V_0\oplus RB)$, что и есть~\eqref{eq:mink}. Замкнутая форма
получается из $\vol(V_0\oplus RB)^{1/n}\ge(\det\Lambda)^{1/n}+R\,\omega_n^{1/n}$.
\end{proof}

Это первая в настоящей работе \emph{строгая} нижняя граница индекса, привязанная
к конкретной родительской решётке: раньше минимальность индекса устанавливалась
только исчерпывающим перебором (например, $1\,424\,115\,231$ подрешёток индекса
$140$ в $A_5^*$). Ценность~\eqref{eq:mink} --- диагностическая. Отношение
достигнутого индекса к правой части показывает, насколько конструкция далека от
объёмного препятствия:

\begin{center}\small
\begin{tabular}{c l r r r}
\toprule
$n$ & решётка & граница~\eqref{eq:mink} & известный $k$ & отношение\\
\midrule
2 & $A_2$    & $4{,}51$   & $7$     & $1{,}55$\\
3 & $A_3^*$  & $10{,}42$  & $15$    & $1{,}44$\\
4 & $D_4$    & $28{,}88$  & $49$    & $1{,}70$\\
5 & $A_5^*$  & $55{,}55$  & $140$   & $2{,}52$\\
6 & $E_6^*$  & $127{,}86$ & $343$   & $2{,}68$\\
7 & $E_7^*$  & $340{,}2$  & $1372$  & $4{,}03$\\
8 & $E_8$    & $663{,}9$  & $2401$  & $3{,}62$\\
9 & $A_9^*$  & $1639{,}7$ & $17253$ & $\mathbf{10{,}52}$\\
\bottomrule
\end{tabular}
\end{center}

\noindent Размерность~$9$ --- очевидный выброс. Вывод, который мы делаем из
этой таблицы: ей не хватало не оптимизации, а \emph{конструкции}. Отметим сразу
и ограничение экрана: он оценивает потенциал, но не достижимость. Так, для
$A_7^*$ граница~\eqref{eq:mink} равна $300{,}7$ против $340{,}2$ у $E_7^*$,
однако наилучший найденный индекс на $A_7^*$ равен $2016$ --- заметно хуже, чем
$1323$ на $E_7^*$.

\subsection{Ламинирование}\label{ssec:lam}

Конструкция строит раскраску $\R^n$ из раскраски $\R^{n-1}$. Пусть
$\Lambda'\subset\R^{n-1}$ --- родительская решётка, $\Gamma'=\ker\psi$ ---
правильная подрешётка индекса~$k$, $c\in\R^{n-1}$ --- \emph{сдвиг слоя},
$t>0$ --- \emph{высота слоя}. Положим
\[
\Lambda=\bigl\{(x+ic,\;it)\ :\ x\in\Lambda',\ i\in\Z\bigr\},\qquad
\Gamma=\bigl\{(x+ic,\;it)\ :\ \psi(x)+ia=0,\ m\mid i\bigr\},
\]
где $a\in\Lambda'/\Gamma'$ --- \emph{склейка}, $m$ --- \emph{модуль слоя};
тогда $[\Lambda:\Gamma]=km$. Как группа $\Lambda=\Lambda'\oplus\Z g$ с
$g=(c,t)$, поэтому $(x,i)\mapsto(\psi(x)+ia,\ i\bmod m)$ действительно
характер.

Работу делают два неравенства, каждое --- в одну строку.

\begin{lemma}\label{lem:P1}
\textbf{[Т]} $\diam\Lambda\le\sqrt{\diam(\Lambda')^2+t^2}$ при любом сдвиге $c$.
\end{lemma}

\begin{proof}
Для точки $(x,z)$ возьмём ближайший слой $i$ (так что $|z-it|\le t/2$) и
ближайшую к $x-ic$ точку решётки $\Lambda'$; тогда
$\dist^2\bigl((x,z),\Lambda\bigr)\le R(\Lambda')^2+t^2/4$, то есть
$R(\Lambda)^2\le R(\Lambda')^2+t^2/4$.
\end{proof}

\begin{lemma}\label{lem:P2}
\textbf{[Т]} $V_0(\Lambda)\subseteq V_0(\Lambda')\times\R$, поэтому для
\emph{горизонтального} вектора $v$ (нулевая слоевая координата, то есть
$v\in\Gamma'$) выполнено $D_\Lambda(v)\ge D_{\Lambda'}(v)$.
\end{lemma}

\begin{proof}
Из $\Lambda'\times\{0\}\subset\Lambda$: если $(x,z)\in V_0(\Lambda)$, то для
всякого $e\in\Lambda'$ имеем $\|x\|^2+z^2\le\|x-e\|^2+z^2$, значит
$x\in V_0(\Lambda')$. Для горизонтального $v$ точка $v/2$ горизонтальна, и
$\dist\bigl(v/2,S\times\R\bigr)=\dist(v/2,S)$.
\end{proof}

Смысл пары лемм: горизонтальные расстояния при подъёме \emph{не убывают}, а
ячейка растёт ровно на $t^2$ по квадрату диаметра. Поэтому запас $d(\Gamma')>1$
базовой раскраски --- это в точности бюджет на рост ячейки, а численно проверять
нужно только векторы с ненулевой слоевой координатой; их единицы, поскольку
$D(v)\ge\|v\|-\diam$ автоматически закрывает всё длиннее $2\diam$.

\subsection{$\chi(\R^9)\le9604$}\label{ssec:dim9}

Базой служит эйзенштейнова раскраска $E_8/2401$ с $d=\sqrt{7/6}$; в
нормировке $\lambda_1(E_8)^2=3$ имеем $\diam E_8=\sqrt6$ и
$D\bigl((3+\omega)E_8\bigr)=\sqrt7$. По лемме~\ref{lem:P1}
$\diam\Lambda\le\sqrt{6+t^2}$, по лемме~\ref{lem:P2} горизонтальные векторы дают
$D\ge\sqrt7$; значит при $t\le1$ строгий диаметр не превосходит $\sqrt7$, и всё
сводится к слоевым векторам.

\begin{keybox}
\textbf{[Ч]} При $m=4$ (индекс $2401\cdot4=9604$), высоте $t=0{,}93$, сдвиге,
смещённом с орбит глубоких дыр $E_8$, и склейке $a=(5,1,4,5)$ выполнено
$D_{\min}=\sqrt7$ и $\diam\le2{,}6200954$, откуда
\[
\chi\bigl(\R^9,[1,\ell]\bigr)\le9604\quad(\ell\le1{,}0097),
\qquad\text{в частности}\quad \chi(\R^9)\le9604 .
\]
Это улучшает оценку $17253$ работы~\cite{ABPR} на $44\%$.
\end{keybox}

Проверка: индекс подтверждён точным целочисленным определителем; перечислением
Финке--Похста получены все $266$ векторов подрешётки нормы $\le2\diam$, из них
$240$ горизонтальных закрыты леммой~\ref{lem:P2}, а $26$ слоевых пересчитаны
\emph{двумя независимыми процедурами проекции} (Дейкстра и SLSQP), согласие ---
до одиннадцатого знака. Статус \textbf{[Ч]}, а не \textbf{[С]}: обе ключевые
оценки суть теоремы, но проверка слоевых векторов ведётся в плавающей точке;
рациональный сертификат уровня~\S\ref{sec:comp} --- ближайшая нерешённая задача.

Существенны две детали. Во-первых, у $E_8$ \emph{две} орбиты глубоких дыр, и
выбор орбиты меняет результат (при индексе $12005$ они дают $1{,}0397$ против
$0{,}8863$), поэтому сдвиг надо хранить явно. Во-вторых, главный рычаг --- сам
сдвиг: смещение с орбит глубоких дыр разрушает <<дважды глубокую дыру>> (точку,
глубокую сразу для двух соседних слоёв) и опускает настоящий диаметр ниже
границы леммы~\ref{lem:P1}.

\subsection{Лестница ширин в $\R^9$}\label{ssec:stair9}

Высоту можно подбирать так, чтобы слоевые векторы \emph{ровно} проходили; тогда
связывает горизонтальный потолок и $d=\sqrt7/\sqrt{6+t^2}$ --- равенство, а не
оценка. Получается лестница (каждая строка --- оценка
$\chi(\R^9,[1,\ell])\le k$ при $\ell<d$):

\begin{center}\small
\begin{tabular}{r r r l}
\toprule
$k$ & $m$ & $t$ & $d$\\
\midrule
$9604$  & $4$  & $0{,}93$ & $1{,}009792$\\
$12005$ & $5$  & $0{,}69$ & $1{,}039662$\\
$14406$ & $6$  & $0{,}55$ & $1{,}053883$\\
$16807$ & $7$  & $0{,}45$ & $1{,}062345$\\
$19208$ & $8$  & $0{,}39$ & $1{,}066688$\\
$24010$ & $10$ & $0{,}33$ & $1{,}070453$\\
$28812$ & $12$ & $0{,}27$ & $1{,}073621$\\
\midrule
--- & $\infty$ & $\to0$ & $\to\sqrt{7/6}=1{,}080123$\\
\bottomrule
\end{tabular}
\end{center}

\noindent Все строки, кроме первой, достигаются уже на сдвиге в глубокую дыру, и
там граница леммы~\ref{lem:P1} \emph{точна} (дважды глубокая дыра существует),
поэтому величины вычислены, а не оценены. Предел лестницы содержателен: при
$m\to\infty$ высота $t\to0$, решётка вырождается в $E_8\times\R$, и ширина
стремится к $\sqrt{7/6}$ --- ровно к ширине самой конструкции $E_8/2401$.
Ламинирование непрерывно интерполирует между рекордом $\R^8$ и $\R^9$.

\subsection{Вещественный множитель и башня для $n\ge10$}\label{ssec:tower}

Для $\Gamma=m\Lambda$ (индекс $m^n$) минимум $D$ достигается на $mu$ с
минимальным~$u$: минимальные векторы всегда релевантны по Вороному, поэтому
ячейка доходит ровно до $\lambda_1/2$ в этом направлении и
$\dist(mu/2,V_0)=(m-1)\lambda_1/2$.

\begin{proposition}\label{prop:mult}
\textbf{[Т]} $D_{\min}(m\Lambda)=(m-1)\lambda_1$, следовательно
$d=(m-1)/\rho$, где $\rho=\diam V_0/\lambda_1$. При $m=3$ это $d=2/\rho$, и
условие правильности --- ровно $\rho\le2$; при $m=2$ потребовалось бы
$\rho\le1$, что невозможно выше размерности~$1$.
\end{proposition}

Правило $d=2/\rho$ сверено с прямым измерением на двенадцати решётках, совпадение
до шести знаков. Подчеркнём, что оно специфично для \emph{вещественных}
множителей: у $3+\omega$ минимум достигается в другом месте, и там работает
тождество $D^2=\tfrac73\lambda_1^2$ из~\S\ref{sec:extra}.

Ламинирование $E_8$ с сохранением $\lambda_1$ удерживает $\rho<2$ во всех
размерностях. Высота при этом фиксирована условием $R^2+t^2\ge1$, и рекурсия
леммы~\ref{lem:P1} даёт $\diam^2_{n+1}\le\tfrac34\diam^2_n+1$ --- неподвижная
точка~$4$, а старт $\diam^2_8=2$ ниже неё, так что $\rho<2$ и $d>1$ при всех
$n\ge8$.

\begin{keybox}
\textbf{[Ч]} Башня ламинаций $E_8$ с раскраской $\Gamma=3\Lambda$ даёт
\[
\chi\bigl(\R^{10},[1,1{,}1795]\bigr)\le3^{10},\quad
\chi\bigl(\R^{11},[1,1{,}1166]\bigr)\le3^{11},\quad
\chi\bigl(\R^{12},[1,1{,}0755]\bigr)\le3^{12}.
\]
\end{keybox}

Неравенством в этой цепочке является только рекурсия для радиуса покрытия:
$\lambda_1=1$ проверяется напрямую, а $D_{\min}=2$ --- исчерпывающим
перечислением всех векторов $3\Lambda$ нормы $\le2\diam$ ($272$, $336$, $432$ и
$4080$ штук при $n=9,\dots,12$), каждый пересчитан двумя процедурами проекции.
Сам индекс $3^n$ --- классический ориентир Лармана--Роджерса; новизна здесь в
явных \emph{ширинах} интервалов и в единообразии конструкции.

Отметим ловушку, на которой мы сначала получили внешне строгую, но неверную
рекурсию: у радиуса покрытия здесь \emph{две} роли, и подставлять в них одну и ту
же оценку нельзя. Высоту $t^2=1-R^2$ надо брать из \emph{нижней} границы $R$
(иначе слоевой вектор оказывается короче единицы и $\lambda_1$ незаметно падает),
а диаметр распространять по \emph{верхней}. Замкнутая форма
$\diam^2_n=4-2(3/4)^{n-8}$ строга лишь до $n=10$, где радиус покрытия
$\Lambda_9$ известен точно.

\subsection{Отрицательные результаты}\label{ssec:dim912neg}

\paragraph{Ниже $9604$ в строгом режиме --- исчерпывающий экран.}
При $t=1$ лемма~\ref{lem:P1} даёт $\diam\le\sqrt7$ \emph{точно}, поэтому
запрещённое множество, построенное против этого диаметра, делает строгой любую
найденную против него раскраску, и искать можно по всем подрешёткам. Множество
построено: $|F|=498\,084$. Общий поиск (Radon/FFT, модули из $\{2,3,5\}$ ---
преобразование размера $7^9$ неподъёмно) при десяти индексах $6480\dots9000$ не
дал ничего; минимум неотделённых векторов равен~$4$ при $6561=3^8$. Если же
зафиксировать четыре характера $(3+\omega)E_8$, остаётся отделить лишь
$F'_a=\{f:\psi(f)+f_9a\equiv0\}$ --- медиана $168$ векторов вместо
$498\,084$, и всё пространство перебирается полностью: $2401$ склейка
$\times\,9841$ проективная форма. Результат: индексы $4802$ и $7203$
\textbf{невозможны}, а $9604$ ($\Z/4$) и $12005$ ($\Z/5$) воспроизводятся, что
служит положительным контролем. Итого \textbf{[Э]}: на этой родительской решётке
$9604$ --- точный пол.

\paragraph{Эйзенштейнов путь в комплексном ранге~$5$.}
Индекс $7^{n/2}$ требует $\rho\le\sqrt{7/3}$. Наилучший тернарный код длины~$5$
с $d\ge3$ имеет $k=2$ (кода $[5,3,3]_3$ не существует по границе Синглтона), и
решётка даёт $\rho=\tfrac23\sqrt7=1{,}7638$; прямое измерение подтверждает
$D_{\min}=\sqrt7$ и $d=0{,}8367$. Вместе с $K_{12}$ (\S\ref{sec:extra}) это
распространяет частичный отрицательный ответ на открытый вопрос~2
работы~\cite{ABPR} с ранга~$6$ на ранг~$5$.

\paragraph{Сертификация радиуса покрытия в $\R^9$.}
Кандидат индекса $7203$ (отношение $1{,}0166$ на измеренном диаметре) остаётся
недоказанным: истинный диаметр зажат между насыщенной LP-оценкой $2{,}602570$ и
строгой границей $2{,}706012$, а порог $\sqrt7=2{,}645751$ лежит строго внутри.
Цену вопроса показывает контрпример из той же кампании: при $t=1{,}30$ бюджет в
$150$ направлений дал отношение $1{,}0016$, а пересчёт с бюджетом $800$ ---
$0{,}9968$. Четыре независимых подхода к сертификации закрыты: ветвление по
направлениям экспоненциально по существу (нужно разрешение $\le21^\circ$ всюду
по $S^8$); перечисление ячеек Делоне требует $\sim10^6$--$10^7$ восхождений
($8000$ восхождений дают $7279$ \emph{различных} вершин, полное число отвечает
$9!=362\,880$); двухслойное усиление бесполезно, так как $r_2^2=1{,}5=R^2(E_8)$
для обеих орбит порядка~$2$; наконец, экран через структуру дыр $E_8$ требует
сдвига, удалённого на $0{,}5$ от всех разностей глубоких дыр, а те плотны с
шагом~$0{,}18$.
